\section{Assignment Problem Definition}

\subsection{Sets, variables, and objective}

Let
\begin{equation}
\mathcal A=\{0,\ldots,n-1\},\qquad
\mathcal R=\{0,\ldots,m-1\}
\end{equation}
be disjoint sets of agents and resources. The binary decision variable
\begin{equation}
x_{ij}=\begin{cases}
1,&\text{if agent }i\text{ is assigned to resource }j,\\
0,&\text{otherwise}
\end{cases}
\end{equation}
is defined for every $(i,j)\in\mathcal A\times\mathcal R$. Given an affinity matrix $S\in\mathbb R^{n\times m}$, the classical objective is
\begin{equation}
\max_x f(x)=\sum_{i=0}^{n-1}\sum_{j=0}^{m-1}S_{ij}x_{ij}.
\label{eq:objective}
\end{equation}

Every supported instance enforces exactly one resource per agent:
\begin{equation}
\sum_{j=0}^{m-1}x_{ij}=1,
\qquad \forall i\in\mathcal A.
\label{eq:agent_constraint}
\end{equation}

\subsection{One-to-one assignment}

For one-to-one matching, every resource is used at most once:
\begin{equation}
\sum_{i=0}^{n-1}x_{ij}\leq 1,
\qquad \forall j\in\mathcal R.
\label{eq:one_to_one_resource}
\end{equation}
When $n=m$, Eqs.~\eqref{eq:agent_constraint}--\eqref{eq:one_to_one_resource} define permutation matrices.

\subsection{Exact-capacity assignment}

For the generalized case, each resource receives a prescribed integer capacity:
\begin{equation}
\bm C=(C_0,\ldots,C_{m-1}),\qquad C_j\in\mathbb Z_{\geq0},
\qquad \sum_{j=0}^{m-1}C_j=n,
\end{equation}
and the resource constraints are
\begin{equation}
\sum_{i=0}^{n-1}x_{ij}=C_j,
\qquad \forall j\in\mathcal R.
\label{eq:capacity_constraint}
\end{equation}
The exact-capacity model is the fixed one-to-$N$ extension used in this phase.

\subsection{Feasible set}

For an instance $I=(S,\bm C,\text{mode})$, let $\mathcal F_I$ denote the set of binary matrices satisfying Eq.~\eqref{eq:agent_constraint} and the applicable resource constraints. The cardinalities of the three fixed feasible sets are
\begin{equation}
|\mathcal F_{3\times3}|=3!=6,
\qquad
|\mathcal F_{4\times4}|=4!=24,
\qquad
|\mathcal F_{4\times2}|=\binom{4}{2}=6.
\end{equation}

\subsection{Benchmark instances}

The fixed matrices are
\begin{align}
S^{(3\times3)}&=
\begin{pmatrix}
0.90&0.20&0.40\\
0.30&0.85&0.50\\
0.45&0.35&0.95
\end{pmatrix},
&\bm C&=(1,1,1),
\label{eq:S33}\\[0.5em]
S^{(4\times4)}&=
\begin{pmatrix}
0.29&0.66&0.52&0.45\\
0.43&0.78&0.87&0.29\\
0.67&0.39&0.92&0.89\\
0.66&0.75&0.56&0.81
\end{pmatrix},
&\bm C&=(1,1,1,1),
\label{eq:S44}\\[0.5em]
S^{(4\times2)}&=
\begin{pmatrix}
0.85&0.30\\
0.40&0.90\\
0.70&0.55\\
0.25&0.80
\end{pmatrix},
&\bm C&=(2,2).
\label{eq:S42}
\end{align}
The $4\times4$ matrix is generated by NumPy's default random generator with seed 2026 and uniform range $[0.15,0.95]$, rounded to two decimals. The $3\times3$ and $4\times2$ matrices are explicit reference instances.

For the capacitated instance, the optimal assignment is
\begin{equation}
X^\star=
\begin{pmatrix}
1&0\\
0&1\\
1&0\\
0&1
\end{pmatrix},
\qquad
f(X^\star)=0.85+0.90+0.70+0.80=3.25.
\label{eq:optimal_assignment_42}
\end{equation}
With the notebook's little-endian row-major convention, this matrix is reported as bitstring \texttt{10011001}.
